# Profil J3 — Consultant senior, 3 ans d'ancienneté

**Fonction :** Consultant senior  
**Ancienneté :** 3 ans  
**Missions :** Cybersécurité pour une institution européenne, documents Word (peu de détails sur les missions par souci de confidentialité)  
**Rapport à l'IA :** Pragmatique — outil très pratique mais nécessitant un encadrement fort pour éviter les hallucinations en livrable client  
**Tic de langage :** "Dans tous les cas"  
**Habitude notable :** Bloc-notes sur son PC avec les meilleurs prompts pour les tâches répétitives  
**Contexte spécifique :** Suit une formation Microsoft sur l'IA pour obtenir une certification et se valoriser sur les missions (forte demande client). A eu des expériences négatives, personnelles et observées chez des collègues, avec des hallucinations et des coquilles dans des livrables clients.  
**Outils utilisés :** Copilot officiel ; Claude officieusement (sans données clients)

---

## T0 — Cadrage (parcours et missions)

> "Je suis consultant senior ici depuis trois ans. Je travaille principalement avec une institution européenne sur des sujets de cybersécurité. C'est un environnement très exigeant — il y a une vraie culture de la confidentialité et de la rigueur. Je ne peux pas entrer dans les détails des missions, mais disons que la plupart de mon travail se fait sur des documents Word, avec des livrables très normés. Dans tous les cas, il y a une procédure à respecter pour chaque rendu."

**Codes fléchés :** (Cadrage — information contextuelle)

---

## T1 — Usage réel (Appropriation)

### Ouverture — dernière fois que vous avez utilisé un outil génératif

> "La semaine dernière, sur un document de synthèse. J'avais une cinquantaine de pages de notes techniques à transformer en un résumé opérationnel pour un comité de direction. J'ai utilisé Copilot — je lui ai filé le document, je lui ai demandé de me sortir une structure avec les grands enjeux et les recommandations. Il m'a proposé une trame. J'ai passé une heure à la retravailler, à recouper avec mes notes, à vérifier que rien n'avait été déformé. Dans tous les cas, je ne fais jamais confiance au premier jet. L'IA est un assistant, pas un rédacteur."

**Codes fléchés :** A1.2 (Synthèse et reformulation), A1.5 (Structuration du raisonnement), A4.1 (Utilité perçue), A4.2 (Facilité d'usage perçue — positive)

### Relance — sur quelles tâches au quotidien ?

> "Je l'utilise pour trois choses principales. La première, c'est la reformulation de textes techniques en langage plus accessible — les clients n'ont pas toujours le niveau technique des équipes. La deuxième, c'est la relecture : je lui demande de détecter des incohérences ou des formulations ambiguës. La troisième, c'est la génération de premiers jets sur des sujets que je maîtrise moins. Dans tous les cas, il y a des sujets où je ne l'utilise jamais : tout ce qui touche à l'architecture de sécurité, aux vulnérabilités spécifiques, ou aux données classifiées. Là, c'est main humaine uniquement."

**Codes fléchés :** A1.2 (Synthèse et reformulation), A1.3 (Première rédaction), A2.1 (Sanctuarisation — confidentialité), A2.2 (Sanctuarisation — jugement)

### Relance — comment as-tu appris à t'en servir ?

> "J'ai commencé en tâtonnant, comme tout le monde. Mais très vite, j'ai compris qu'il fallait une approche structurée. J'ai ouvert un bloc-notes sur mon PC où je note les prompts qui fonctionnent bien pour les tâches récurrentes — par exemple, un prompt pour synthétiser des notes techniques, un autre pour reformuler en langage client, un autre pour détecter des incohérences. Je les affine au fil du temps. Et depuis quelques mois, je suis une formation certifiante Microsoft sur l'IA. Pas seulement pour apprendre, mais aussi pour avoir un papier à valoriser. La demande des clients est de plus en plus forte sur ces sujets, donc c'est un bon investissement."

**Codes fléchés :** A3.1 (Apprentissage autonome), A3.3 (Formation institutionnelle), E2.1 (Prospective du métier — il anticipe la demande client)

### Relance — le temps gagné, tu l'as mis sur quoi ?

> "Sur le document de synthèse dont je parlais, j'ai gagné environ deux heures sur la première structuration. Je les ai utilisées pour faire une relecture croisée avec les sources originales. Parce que l'IA peut faire des erreurs subtiles — des formulations qui semblent correctes mais qui trahissent le sens. Dans tous les cas, je préfère passer plus de temps à vérifier qu'à produire. La qualité prime sur la vitesse."

**Codes fléchés :** A4.1 (Utilité perçue), B2.3 (Vérification comme compétence), D2.2 (Vérification des sources)

### Relance — une situation où ça t'a demandé plus d'effort que de le faire toi-même ?

> "Oui, plusieurs fois. La plus marquante, c'était sur un document où j'avais demandé à Copilot de reformuler un passage technique. Il m'a sorti une version qui sonnait bien, mais en relisant, je me suis rendu compte qu'il avait inversé deux concepts clés. Le sens était complètement faussé. Si j'avais envoyé ça au client, ça aurait été une catastrophe. J'ai passé quarante-cinq minutes à tout reprendre, alors que j'aurais pu le faire moi-même en vingt minutes. Depuis, je suis encore plus vigilant."

**Codes fléchés :** A4.2 (Facilité d'usage perçue — négative), D3.1 (Hallucination détectée), A4.3 (Résistance ou évitement — il devient plus vigilant)

### Relance — vous en parlez entre vous ?

> "Entre consultants, oui, mais pas dans les grandes largeurs. On se partage des retours d'expérience, des mises en garde. J'ai déjà alerté des collègues juniors sur le fait qu'il ne faut pas faire confiance aveuglément. J'en ai vu, des livrables avec des coquilles grossières parce que le consultant n'avait pas relu ce que l'IA avait produit. Dans tous les cas, je dis toujours : 'si tu ne peux pas justifier ce que l'IA a écrit, ne le mets pas dans le livrable'. C'est une règle que j'essaie de transmettre."

**Codes fléchés :** A3.2 (Transmission entre pairs), E1.1 (Écart discours-pratique — les juniors ne relisent pas assez), C2.2 (Opacité stratégique — on n'en parle pas aux clients)

---

## T2 — Apprentissage du métier (Compétences)

### Ouverture — comment on apprend le métier aujourd'hui ?

> "Le métier, on l'apprend surtout par l'expérience et par les erreurs. Moi, mes premières missions, j'ai fait des erreurs, j'ai appris. L'IA change un peu la donne : elle te donne une première version tellement propre que tu peux avoir l'illusion de maîtriser un sujet sans l'avoir vraiment creusé. Dans tous les cas, un bon consultant, c'est quelqu'un qui sait ce qu'il met dans ses livrables, pas quelqu'un qui sait formater un prompt. La formation continue — et celle que je suis actuellement — me conforte dans cette idée : l'IA est un outil, pas une compétence en soi."

**Codes fléchés :** B1.1 (Vertu formatrice reconnue — l'expérience terrain reste clé), B2.2 (Atrophie du raisonnement — émergeant, il craint l'illusion de maîtrise), B2.1 (Accélération de la montée en compétence — ambivalent)

### Relance — par quelles tâches on apprend réellement le métier ?

> "Celles où tu es seul face au problème. Quand tu dois produire une analyse, justifier une recommandation, expliquer une vulnérabilité. Là, l'IA ne peut pas t'aider vraiment, parce qu'elle ne connaît pas le contexte spécifique du client, ses contraintes, ses sensibilités. C'est dans ces moments que tu progresses. Les tâches répétitives — reformulation, mise en forme — elles ne sont pas formatrices. L'IA les fait très bien, et tant mieux. Mais le cœur du métier, c'est le jugement."

**Codes fléchés :** B1.1 (Vertu formatrice reconnue — pour les tâches complexes), B1.2 (Vertu formatrice déniée — pour les tâches répétitives), A2.2 (Sanctuarisation — jugement)

### Relance — les tâches ingrates, sont-elles encore formatrices ?

> "Pour les juniors, oui. Passer des heures à structurer un document, à chercher la bonne formulation, à vérifier chaque source — c'est comme ça qu'on développe le réflexe de la rigueur. Mais moi, avec trois ans d'expérience, j'ai déjà ce réflexe. Donc je n'hésite pas à déléguer la partie rédactionnelle à l'IA pour me concentrer sur le fond. Dans tous les cas, je vérifie tout. Mais je ne perds plus mon temps sur ce que l'outil peut faire aussi bien."

**Codes fléchés :** B1.1 (Vertu formatrice reconnue — pour les juniors, pas pour lui), B2.3 (Vérification comme compétence — il a acquis le réflexe)

### Relance — comment vous vérifiez ce que produit l'outil ?

> "Je le fais de plusieurs façons. D'abord, une relecture de cohérence : est-ce que le texte tient debout ? Ensuite, un contrôle croisé avec les sources : est-ce que les informations sont exactes ? Et enfin, je me pose une question : 'est-ce que j'oserais défendre ce passage en réunion avec le client ?' Si la réponse est non, je le reprends. Dans tous les cas, la règle est simple : si je ne peux pas l'expliquer, je ne le laisse pas dans le livrable."

**Codes fléchés :** B2.3 (Vérification comme compétence), D2.2 (Vérification des sources), D4.1 (Attribution de responsabilité — il sait que c'est sa responsabilité)

---

## T3 — Client et valeur (Légitimité de l'expert)

### Ouverture — la confiance du client repose sur quoi ?

> "La confiance repose sur la fiabilité. Dans la cybersécurité, si tu te trompes, les conséquences sont graves. Le client doit pouvoir se dire : 'ces gens-là savent ce qu'ils font, ils ne prennent pas de risques inutiles'. L'IA peut aider à produire plus vite, mais elle ne peut pas remplacer le jugement d'un expert. La confiance, elle vient de la capacité à expliquer, à justifier, à prendre position."

**Codes fléchés :** C1.1 (Légitimité cognitive — l'expertise technique), C1.3 (Légitimité morale — la responsabilité)

### Relance — diriez-vous à un client qu'une partie a été produite par l'IA ?

> "Non, et je ne le ferai pas. Non pas parce que je veux cacher quelque chose, mais parce que je pense que c'est contre-productif. Le client paie une expertise, pas un outil. Si je lui dis 'c'est l'IA qui a fait ça', il va se demander pourquoi il nous paie. Dans tous les cas, le produit final est validé par nous, donc c'est notre responsabilité. La question de l'origine du texte est secondaire."

**Codes fléchés :** C2.2 (Opacité stratégique), C1.3 (Légitimité morale — la responsabilité du produit final)

### Relance — si un livrable était perçu comme moins crédible, pourquoi ?

> "Cela viendrait d'une erreur factuelle. Une donnée fausse, une référence inexacte, un concept mal interprété. Là, le client perd confiance, et il a raison. Dans la cybersécurité, la moindre erreur peut être exploitable. Donc je suis extrêmement vigilant. Et d'ailleurs, c'est pour ça que j'ai vu des livrables avec des coquilles — des collègues qui n'avaient pas relu ce que l'IA avait écrit. Dans tous les cas, un livrable mal vérifié, c'est une faute professionnelle, que le texte vienne de l'IA ou de la main du consultant."

**Codes fléchés :** C1.1 (Légitimité cognitive — l'exactitude est clé), D3.1 (Hallucination détectée — écho aux expériences vécues), D4.1 (Attribution de responsabilité)

### Relance — dans dix ans, sur quoi un client acceptera-t-il de payer un cabinet ?

> "Il paiera pour le jugement et l'engagement. L'IA pourra produire des rapports, faire des analyses, proposer des scénarios. Mais elle ne pourra pas prendre de décision à la place d'un expert, ni assumer la responsabilité d'une recommandation. Dans mon secteur, on a besoin de gens qui disent 'voilà ce qu'il faut faire, et voilà pourquoi'. Je ne vois pas une IA faire ça un jour."

**Codes fléchés :** C3.1 (Déplacement vers le jugement), C3.2 (Déplacement vers la responsabilité), E2.1 (Prospective du métier)

---

## T4 — Gouvernance et risques

### Ouverture — il existe des règles sur l'usage des outils ici ?

> "Oui, il y a une politique officielle. Copilot est approuvé pour un usage général, à condition de ne pas y mettre de données sensibles ou classifiées. Mais pour être honnête, la règle est assez générale. Dans mon contexte, avec une institution européenne, les exigences de confidentialité sont encore plus fortes. Donc j'applique une règle plus stricte que ce que le cabinet demande. Dans tous les cas, je préfère la prudence."

**Codes fléchés :** D1.1 (Règle formelle connue), D3.3 (Enjeu de confidentialité), D1.3 (Flou normatif — la règle est générale)

### Relance — avant qu'un livrable ne parte chez le client, il se passe quoi ?

> "Il y a un circuit de validation. En général, je produis une version, je la soumets à un relecteur — souvent un manager ou un expert technique — puis on valide ensemble. Mais il n'y a pas de procédure spécifique pour vérifier les contenus générés par l'IA. C'est au consultant de s'assurer que tout est correct. Dans tous les cas, c'est moi qui signe le livrable, donc c'est moi qui suis responsable."

**Codes fléchés :** D2.1 (Revue hiérarchique), D2.3 (Absence de contrôle dédié), D4.1 (Attribution de responsabilité)

### Relance — déjà rencontré un contenu produit par l'IA qui s'est révélé faux ?

"Oui, plusieurs fois. La plus récente, c'était une synthèse de Copilot où il avait mal interprété un point technique. Rien de dramatique — je l'ai repéré à la relecture. Mais j'ai aussi vu des collègues juniors se faire avoir. Ils envoient un livrable, le manager voit une incohérence, il demande des comptes. Et là, le junior répond 'c'est l'IA qui a fait'. Ça n'a jamais bien fini. On est formateur, pas un journal de bord."

**Codes fléchés :** D3.1 (Hallucination détectée — vécue et observée), D3.2 (Quasi-incident — ça n'est pas parti au client), D4.1 (Attribution de responsabilité — les juniors qui se cachent derrière l'IA)

### Relance — si vous écriviez la règle vous-même, vous diriez quoi ?

"Je dirais : 'L'IA est un outil d'assistance, pas un substitut à l'expertise. Toute production doit être vérifiée et justifiée par le consultant. En cas de doute, l'IA n'est pas utilisée.' Et j'ajouterais une obligation de traçabilité : il faut pouvoir dire ce qui a été fait par l'IA et ce qui ne l'a pas été. Pas pour contrôler, mais pour responsabiliser."

**Codes fléchés :** D5.2 (Demande de cadrage), E2.2 (Recommandation spontanée), D2.2 (Vérification des sources)

---

## T5 — Clôture

### Relance — à la place de la direction, quelle serait votre première décision ?

"Je mettrais en place une formation obligatoire, pas sur l'utilisation technique, mais sur la détection des erreurs. Parce que le vrai risque, ce n'est pas que l'IA fasse des erreurs — c'est que les consultants ne les voient pas. Et je formaliserais un processus de validation des livrables qui inclut une vérification spécifique sur les contenus générés par l'IA. Un peu comme on fait pour les citations, mais en plus systématique."

**Codes fléchés :** E2.2 (Recommandation spontanée), D5.2 (Demande de cadrage)

### Relance — un aspect important qu'on n'a pas abordé ?

"La dimension juridique. Avec l'IA, il y a des questions de propriété intellectuelle, de confidentialité, de responsabilité. Personne n'en parle vraiment, mais c'est un sujet qui va devenir central, surtout dans mon secteur. Dans tous les cas, il y a un vide juridique que le cabinet va devoir combler — soit en anticipant, soit en subissant."

**Codes fléchés :** E2.1 (Prospective du métier), D3.3 (Enjeu de confidentialité), E1.1 (Écart discours-pratique — le sujet n'est pas abordé)

### Relance — quelqu'un dont le point de vue serait particulièrement utile ?

"Un responsable juridique du cabinet. Pour voir comment il perçoit les risques — confidentialité, propriété intellectuelle, responsabilité en cas d'erreur. Parce que moi, je vois les risques opérationnels, mais je ne vois pas les risques légaux. Et je pense que ces risques sont beaucoup plus grands que ce qu'on imagine."

**Codes fléchés :** E2.1 (Prospective du métier), D3.3 (Enjeu de confidentialité), D4.1 (Attribution de responsabilité — sous-jacent)